\section{Projective flags and selected-stream entropy}
\label{sec:flags}

We first construct the graph and record the entropy that every long
consistent tuple extracted from its random stream must retain.
This lower bound will be compared with descriptions produced later.

\subsection{Incidence estimates}

For a prime power \(q\), let \(\mathcal P_d=\PG(d,q)\) be the
one-dimensional subspaces of \(\FF_q^{d+1}\), and let
\(\mathcal P_d^*\) be its dual point space. A dual point \(b\)
represents the hyperplane \(b^\perp\). We write \(a\perp b\)
for incidence. In both spaces, a projective \(j\)-flat corresponds
to a vector subspace of dimension \(j+1\).
Put \(Q_j=1+q+\cdots+q^j\).

\begin{lemma}[Incidence and variance]
\label{lem:incidence}
Let \(d\ge2\), and put \(p_d=Q_{d-1}/Q_d=(1-1/Q_d)/q\).
For real weights \(w_a\) on \(\mathcal P_d\), set
\(W(b)=\sum_{a\perp b}w_a\) and \(w=\sum_a w_a\). Then
\[
 \frac1{Q_d}\sum_b W(b)=p_dw,\qquad
 \sum_b\bigl(W(b)-p_dw\bigr)^2
 \le q^{d-1}\sum_a w_a^2.
\]
The same assertions hold with the two spaces interchanged.
In particular, for \(A\subseteq\mathcal P_d\) and
\(B\subseteq\mathcal P_d^*\),
\[
 \left|e(A,B)-p_d|A||B|\right|
 \le q^{(d-1)/2}\sqrt{|A||B|},
 \qquad
 e(A,B)=|\{(a,b)\in A\times B:a\perp b\}|.
\]
If \(A,B\ne\varnothing\) and their incidence density is \(o(1/q)\),
then \(|A||B|\le Cq^{d+1}\) for all sufficiently large \(q\).
\end{lemma}

\begin{proof}
Every point lies on \(Q_{d-1}\) hyperplanes and every two distinct
points lie on \(Q_{d-2}\) hyperplanes. Thus the incidence matrix
\(M_d\) satisfies
\[
 M_dM_d^{\mathsf T}=q^{d-1}I+Q_{d-2}\mathbf1\mathbf1^{\mathsf T}.
\]
Both row and column sums equal \(Q_{d-1}\). Subtracting the constant
part of a weight vector therefore gives the variance estimate.
Applying this operator bound to the centered indicator vectors of
\(A\) and \(B\) gives the displayed mixing estimate. Finally,
\(p_d\ge1/(2q)\), so a density \(o(1/q)\) forces
\(\sqrt{|A||B|}\le Cq^{(d+1)/2}\).
\end{proof}

This projective incidence calculation is the bipartite form of the
orthogonality-graph spectral calculation of Alon and Krivelevich
\cite[Section~2]{AlonKrivelevich1997}; it also underlies
Brada\v{c}'s polarity-graph estimates \cite[Lemma~2.1]{Bradac2026}.
We shall use both its unweighted and weighted forms.

\subsection{The ordered graph}

Fix \(0<\eta<1/10\). Throughout the lower-bound proof let \(q\)
be a sufficiently large prime, and put
\[
 \beta=\eta/10^7,\qquad \sigma=\log q,\qquad
 N=\lfloor q^4\sigma\rfloor,\qquad
 k=\lfloor q\sigma^{1+\eta}\rfloor.
\]
The small parameter \(\beta\) is unrelated to the independence number
\(\alpha(G)\). Constants may depend on \(\eta\); they never depend
on \(q\). Each \(o(1)\) is uniform in the stated parameter ranges.
Only a bounded number of stages, depending on \(\eta\), will be used.

A \emph{flag} is a pair \((a,b)\in\mathcal P_4\times\mathcal P_4^*\)
with \(a\perp b\). There are \(M=Q_4Q_3\) flags.
Choose a stream \(X_1,\ldots,X_N\) of independent uniform flags.
Its graph has vertex set \(\{1,\ldots,N\}\); writing
\(X_i=(a_i,b_i)\), its edge rule for \(i<j\) is
\[
 i\sim j
 \quad\Longleftrightarrow\quad
 a_i\perp b_j\ \hbox{ and }\ a_j\not\perp b_i.
\]
Repeated flags are permitted: the vertices are stream positions.
This is the ordered version of Brada\v{c}'s \(D^*\) construction
\cite[Section~2.5]{Bradac2026}.
Related earlier constructions use ordered edges of a bipartite
incidence graph \cite{CodenottiPudlakResta2000,KostochkaPudlakRodl2010};
the reverse nonincidence in our displayed edge rule is part of the
specific construction used here.

\begin{lemma}
\label{lem:clique-free}
Every stream graph is \(K_5\)-free.
\end{lemma}
\begin{proof}
In a clique with increasing positions \(i_1<\cdots<i_5\), the
hyperplane \(b_{i_j}^\perp\) contains
\(a_{i_1},\ldots,a_{i_j}\) but not \(a_{i_{j+1}}\), for
\(j=1,\ldots,4\).
Thus these five points span the full five-dimensional vector space.
They are all annihilated by \(b_{i_5}\), a contradiction.
\end{proof}

An ordered tuple of flags is \emph{consistent} if
\[
 i<j,\ a_i\perp b_j \quad\Longrightarrow\quad a_j\perp b_i.
\]
Independent vertex sets give consistent tuples in stream order.
Reversing the tuple and interchanging its two endpoint systems
preserves consistency.

\begin{lemma}[Rectangle occupancy]
\label{lem:rectangles}
With probability \(1-o(1)\), the following holds simultaneously
for every vector subspace \(V\subseteq(\FF_q^5)^*\) in the dual space:
\[
 \#\{i:X_i\in\mathcal R(V)\}\le C\sigma,\qquad
 \mathcal R(V)=\{(a,b):b\subseteq V,\ a\subseteq V^\perp\}.
\]
The same bound holds for the reverse-swapped stream.
\end{lemma}
\begin{proof}
For \(1\le r=\dim V\le4\), the rectangle contains
\(Q_{r-1}Q_{4-r}\le4q^3\) flags; for \(r=0,5\) it is empty.
Its stream count is binomial with mean at most \(C_0\sigma\).
The number of vector subspaces is at most \(6q^{25}\), for example
by counting ordered spanning lists. A Chernoff bound and a sufficiently
large fixed \(C\) give the simultaneous claim.
After swapping, \(\mathcal R(V)\) becomes
\(\mathcal R(V^\perp)\), so the same event suffices.
\end{proof}

\subsection{Laws of selected tuples}

Entropy always means Shannon entropy with natural logarithms.
For a finite-valued variable \(Y\) with masses \(p_y\), write
\(H(Y)=-\sum_y p_y\log p_y\), with \(0\log0=0\).
Conditional entropy \(H(Y\mid Z)\) averages this quantity over \(Z\);
\(H_\theta(Y)\) denotes the entropy of the particular conditional law
given \(Z=\theta\). Mutual information and total correlation are
\[
 I(Y;Z)=H(Y)+H(Z)-H(Y,Z),\qquad
 \TC(Y_1,\ldots,Y_m)=\sum_iH(Y_i)-H(Y_1,\ldots,Y_m).
\]
We use their conditional versions with the same convention.

For probability masses \(P,Q\) on a finite set, their relative entropy is
\[
 \KL(P\Vert Q)=\sum_xP(x)\log\frac{P(x)}{Q(x)},
\]
with value infinity if \(P\) charges a zero of \(Q\). Direct expansion gives
\[
 I(Y;Z)=\KL(P_{Y,Z}\Vert P_Y\otimes P_Z).
\]
For every real function \(f\), Jensen's inequality gives
\begin{equation}\label{eq:entropy-variational}
 \KL(P\Vert Q)\ge \EE_P f-\log\EE_Q e^f.
\end{equation}
Indeed, apply concavity of the logarithm under \(P\) to
\(e^fQ/P\); its expectation is at most \(\EE_Q e^f\).
In particular relative entropy is nonnegative. If \(P(E)=0\),
use \(f=-s1_E\) and let \(s\to\infty\) to obtain
\(\KL(P\Vert Q)\ge-\log(1-Q(E))\ge Q(E)\).

We also use the following elementary concentration estimates.
For a binomial or Poisson variable of mean \(\mu\), halving or
doubling its mean has tail probability at most \(e^{-c\mu}\);
more generally, its upper tail at \(h\ge e\mu\) is at most
\((e\mu/h)^h\). For independent centered variables bounded in
absolute value by one and with total variance \(V\), the two-sided
tail of their sum at \(z\) is at most
\[
 2\exp\{-c\min(z,z^2/V)\}.
\]
These follow from the exponential Markov inequality. For the first
two bounds use
\(\log\EE e^{t(X-\mu)}\le\mu(e^t-1-t)\).
For the last, the Taylor expansion gives
\(\EE e^{tX}\le\exp(Ct^2\EE X^2)\) for centered \(|X|\le1\)
and \(|t|\le1\); multiply these bounds and choose
\(|t|=\min(cz/V,c)\). The zero-variance case is deterministic.

Say a stream law has \emph{density at most \(C_0\)} if its probability
of every stream event is at most \(C_0\) times the iid uniform
stream probability. Auxiliary random choices can have arbitrary
dependence on the stream; only this marginal bound is required.

\begin{lemma}[Entropy of an extracted tuple]
\label{lem:stream-entropy}
Suppose the stream law has density at most \(C_0\), where \(C_0\)
is independent of \(q\). By any rule, possibly using auxiliary
randomness, extract a tuple \(F\) of deterministic length \(\ell\)
in stream order or in reverse-swapped order. If
\(c_0k\le\ell\le C_1k\) for fixed positive constants, then
\[
 \max_f\PP(F=f)\le
 \frac{2C_0\binom N\ell}{M^\ell},
 \qquad
 H(F)\ge4\sigma\ell+\eta\ell\log\sigma-O(k).
\]
\end{lemma}
\begin{proof}
The event \(F=f\) requires \(f\), or its reverse-swap, to occur at
some \(\ell\) positions of the stream. For any specified positions
the iid probability is \(M^{-\ell}\), even when flag values repeat.
A union bound proves the atom bound without any independence
assumption on the selected tuple. Entropy is at least minus the
logarithm of its largest atom. Since
\(\binom N\ell\le(eN/\ell)^\ell\),
\(\log M=7\sigma+O(1/q)\), and
\(\log\ell=\sigma+(1+\eta)\log\sigma+O(1)\), this gives the result.
\end{proof}

Conditioning on an event of probability at least a fixed positive
constant preserves stream domination, with a different constant.
Fixing independent public random tables by averaging will be done
before such conditioning. A fixed finite number of these operations
therefore leaves Lemma~\ref{lem:stream-entropy} applicable.

To prove the lower bound, suppose every stream has a consistent
\(k\)-tuple, select one deterministically, and condition on the event
of Lemma~\ref{lem:rectangles}. The compression argument starts from
this law. Every subsequent retained tuple will have deterministic
length \(\Theta(k)\), satisfy the same rectangle bound, and have a
stream law with bounded density. Its entropy must consequently
retain the excess \(\eta\ell\log\sigma\).
